% =============================================================================
% SAMVADA: IEEE Research Paper — LaTeX Source
% Compile with: pdflatex -> pdflatex (two passes for references)
% =============================================================================

\documentclass[conference]{IEEEtran}

% -- Packages ----------------------------------------------------------------
\usepackage[utf8]{inputenc}
\usepackage[T1]{fontenc}
\usepackage{cite}
\usepackage{amsmath,amssymb,amsfonts}
\usepackage{graphicx}
\usepackage{textcomp}
\usepackage{xcolor}
\usepackage{booktabs}
\usepackage{array}
\usepackage{multirow}
\usepackage{hyperref}
\usepackage{listings}
\usepackage{float}
\usepackage{tabularx}
\usepackage{url}
\usepackage{fancyvrb}
\usepackage{enumitem}
\usepackage{balance}
\usepackage{tikz}
\usetikzlibrary{shapes.geometric, arrows.meta, positioning, fit, backgrounds}

% -- Listings style for code --------------------------------------------------
\lstset{
  basicstyle=\ttfamily\scriptsize,
  breaklines=true,
  frame=single,
  backgroundcolor=\color{gray!8},
  keywordstyle=\color{blue!70!black},
  commentstyle=\color{green!50!black},
  stringstyle=\color{red!60!black},
  showstringspaces=false,
  tabsize=2,
  captionpos=b,
}

% -- Hyperref config ----------------------------------------------------------
\hypersetup{
  colorlinks=true,
  linkcolor=blue!60!black,
  citecolor=blue!60!black,
  urlcolor=blue!60!black,
}

% -- Custom column types ------------------------------------------------------
\newcolumntype{L}[1]{>{\raggedright\arraybackslash}p{#1}}
\newcolumntype{C}[1]{>{\centering\arraybackslash}p{#1}}
\newcolumntype{R}[1]{>{\raggedleft\arraybackslash}p{#1}}

% =============================================================================
\begin{document}

% -- Title --------------------------------------------------------------------
\title{SAMVADA: A Real-Time Speech-to-Speech Translation Platform for Multilingual Video Conferencing}

% -- Authors ------------------------------------------------------------------
\author{
  \IEEEauthorblockN{Author Name\textsuperscript{1}}
  \IEEEauthorblockA{
    \textsuperscript{1}Department of Computer Science and Engineering\\
    Institution Name, City, Country\\
    Email: author@institution.edu
  }
}

\maketitle

% =============================================================================
% ABSTRACT
% =============================================================================
\begin{abstract}
The proliferation of globalized collaboration demands seamless cross-lingual communication in real-time conferencing environments. This paper presents \textbf{SAMVADA}, an end-to-end, real-time Speech-to-Speech Translation (S2ST) platform that enables participants in a video conference to speak in their native language while receiving translated subtitles and synthesized speech in their preferred language. The system implements a cascaded AI pipeline comprising OpenAI Whisper (Automatic Speech Recognition), Meta NLLB-200 (Neural Machine Translation), and a multi-engine Text-to-Speech synthesis stack (Microsoft Edge Neural TTS, Sarvam AI, Google TTS). The platform supports \textbf{30 languages} including \textbf{12 Indic languages} and achieves sub-4-second end-to-end latency on commodity GPU hardware (NVIDIA T4). Key contributions include: (1)~a novel server-side translation orchestrator with per-speaker priority queuing, (2)~a persistent multiplexed WebSocket architecture eliminating per-utterance connection overhead, (3)~a multi-stage hallucination filtering framework for Whisper ASR, (4)~an AudioWorklet-based Voice Activity Detection (VAD) system with continuous capture, and (5)~a comprehensive pipeline diagnostic engine for real-time quality monitoring. Experimental evaluation on a 10-sentence English-to-Hindi benchmark yields a BLEU score of approximately 25--30 using the NLLB-200-distilled-600M model with CTranslate2 INT8 quantization, while maintaining an average per-sentence NMT latency of under 150\,ms on GPU. The system is deployed as a three-tier web application using React, Node.js, and FastAPI, with MongoDB for persistent storage and Socket.IO for real-time bi-directional communication.
\end{abstract}

% -- Keywords -----------------------------------------------------------------
\begin{IEEEkeywords}
Speech-to-Speech Translation, Multilingual Conferencing, Automatic Speech Recognition, Neural Machine Translation, Text-to-Speech Synthesis, Real-Time Systems, WebSocket, Voice Activity Detection, Whisper, NLLB-200
\end{IEEEkeywords}

% =============================================================================
% I. INTRODUCTION
% =============================================================================
\section{Introduction}

The modern workforce is increasingly distributed across geographies, cultures, and linguistic boundaries. Video conferencing platforms such as Zoom, Microsoft Teams, and Google Meet have become indispensable tools for remote collaboration. However, these platforms predominantly operate in a monolingual paradigm, creating significant barriers for participants who are not fluent in the meeting's primary language. According to Ethnologue~\cite{ethnologue}, there are over 7,000 living languages, yet mainstream conferencing tools offer limited or no real-time translation support beyond post-hoc captioning.

India, with its 22 officially recognized scheduled languages and hundreds of dialects, exemplifies this challenge acutely. A typical corporate meeting in India may involve participants speaking Hindi, Tamil, Telugu, Marathi, Bengali, and English---often within the same conversation. Existing solutions such as Google Translate provide text-based translation but lack the real-time, speech-integrated experience required for fluid multi-party conversations.

The goal of this work is to design and implement \textbf{SAMVADA} (Sanskrit: \textit{Samvada}, meaning ``dialogue'' or ``conversation''), a production-grade multilingual video conferencing platform that provides:

\begin{enumerate}
  \item \textbf{Real-time Automatic Speech Recognition (ASR)} with automatic language detection.
  \item \textbf{Neural Machine Translation (NMT)} to $N$ target languages simultaneously.
  \item \textbf{Neural Text-to-Speech (TTS) synthesis} delivering translated audio in natural human voices.
  \item \textbf{Live translated subtitles} overlaid on the video conferencing interface.
  \item \textbf{Post-meeting AI-generated summaries} using Large Language Models.
\end{enumerate}

The remainder of this paper is organized as follows: Section~\ref{sec:related} reviews related work in speech translation and multilingual conferencing. Section~\ref{sec:architecture} describes the system architecture and design principles. Section~\ref{sec:pipeline} details the AI pipeline components. Section~\ref{sec:implementation} presents implementation specifics across the three-tier application stack. Section~\ref{sec:evaluation} discusses the experimental setup and performance evaluation. Section~\ref{sec:results} analyzes the results and discusses limitations. Section~\ref{sec:conclusion} concludes the paper and outlines future work.

% =============================================================================
% II. RELATED WORK
% =============================================================================
\section{Related Work}
\label{sec:related}

\subsection{Speech-to-Speech Translation Systems}

The field of Speech-to-Speech Translation has evolved from early cascaded systems to modern end-to-end approaches. Traditional cascaded S2ST systems decompose the problem into three sequential stages: ASR $\rightarrow$ MT $\rightarrow$ TTS~\cite{nakamura2006atr}. Recent end-to-end models such as Meta's SeamlessM4T~\cite{seamlessm4t} and Google's Translatotron~2~\cite{translatotron2} attempt to bypass the intermediate text representation, directly mapping source speech to target speech.

However, cascaded architectures remain dominant in production systems due to several advantages: (a)~modular upgradability---each component can be independently improved; (b)~interpretability---intermediate text outputs enable debugging, transcript storage, and subtitle generation; and (c)~flexibility---different TTS engines can be selected per language. SAMVADA adopts the cascaded approach for these reasons while leveraging state-of-the-art models at each stage.

\subsection{Automatic Speech Recognition}

OpenAI's Whisper~\cite{whisper} is a transformer-based ASR model trained on 680,000 hours of multilingual and multitask supervised data. Whisper supports 99~languages and provides both transcription and language identification. The \texttt{faster-whisper} library~\cite{fasterwhisper} re-implements Whisper using CTranslate2, achieving $4\times$ faster inference with comparable accuracy through INT8 quantization. SAMVADA employs faster-whisper with the \texttt{small} model variant, balancing accuracy and latency for real-time applications.

\subsection{Neural Machine Translation}

Meta's No Language Left Behind (NLLB-200)~\cite{nllb} is a family of multilingual NMT models supporting 200+ languages. The distilled 600M-parameter variant offers a practical trade-off between translation quality and inference speed. By converting the model to CTranslate2 INT8 format, SAMVADA achieves 2--4$\times$ speedup over the HuggingFace Transformers baseline while supporting batched translation to multiple target languages simultaneously.

\subsection{Text-to-Speech Synthesis}

Modern neural TTS systems have achieved near-human naturalness. Microsoft's Edge Neural TTS~\cite{edgetts} provides studio-grade voices across 75+ languages with zero cost. Sarvam AI's Bulbul~v3 model~\cite{sarvam} is specifically designed for Indian regional languages, offering superior prosody for Hindi, Tamil, Telugu, and other Indic languages. Google's gTTS~\cite{gtts} serves as a universal fallback. SAMVADA implements a multi-engine TTS architecture that selects the optimal engine based on target language and availability.

\subsection{Real-Time Conferencing Platforms}

WebRTC~\cite{webrtc} is the standard for peer-to-peer audio/video communication in browsers. Socket.IO~\cite{socketio} provides reliable real-time bi-directional event-based communication over WebSockets. Existing multilingual conferencing solutions like Wordly.ai and KUDO offer interpretation services but typically require human interpreters or operate with significant latency. SAMVADA differentiates itself by providing fully automated, sub-4-second S2ST with no human intervention.

% =============================================================================
% III. SYSTEM ARCHITECTURE
% =============================================================================
\section{System Architecture}
\label{sec:architecture}

\subsection{Architectural Overview}

SAMVADA follows a \textbf{three-tier microservices architecture} consisting of:

\begin{enumerate}
  \item \textbf{Client Tier}---React-based Single Page Application (SPA).
  \item \textbf{Server Tier}---Node.js/Express backend with Socket.IO.
  \item \textbf{AI Service Tier}---Python FastAPI server hosting the ML pipeline.
\end{enumerate}

Fig.~\ref{fig:architecture} illustrates the high-level data flow. The client captures audio via an AudioWorklet processor, encodes it as WAV, and transmits it to the server via Socket.IO binary events. The server's Translation Orchestrator computes the required target languages, forwards the audio to the FastAPI AI service over a persistent WebSocket, and routes the translated text and synthesized audio back to the appropriate participants.

\begin{figure}[htbp]
\centering
\resizebox{0.95\columnwidth}{!}{%
\begin{tikzpicture}[
  tier/.style={rectangle, rounded corners=4pt, draw=#1!60, fill=#1!8,
               minimum width=6.8cm, minimum height=2.2cm, align=center,
               font=\small, line width=0.8pt},
  arr/.style={-{Stealth[length=5pt]}, thick, draw=gray!70},
  lbl/.style={font=\scriptsize\sffamily, text=gray!80},
]
  % Client Tier
  \node[tier=blue] (client) {
    \textbf{CLIENT TIER}\\
    \scriptsize React 19 $\cdot$ Vite 8 $\cdot$ Socket.IO Client\\
    \scriptsize AudioWorklet VAD $\cdot$ SubtitleOverlay
  };
  % Server Tier
  \node[tier=teal, below=1.2cm of client] (server) {
    \textbf{SERVER TIER}\\
    \scriptsize Node.js $\cdot$ Express $\cdot$ Socket.IO $\cdot$ MongoDB\\
    \scriptsize Translation Orchestrator (Per-Speaker Queue)
  };
  % AI Tier
  \node[tier=violet, below=1.2cm of server] (ai) {
    \textbf{AI SERVICE TIER}\\
    \scriptsize FastAPI $\cdot$ PyTorch $\cdot$ CUDA\\
    \scriptsize Whisper $\rightarrow$ NLLB-200 $\rightarrow$ Edge TTS / Sarvam AI
  };
  % Arrows
  \draw[arr] (client.south) -- node[lbl, right=2pt] {Socket.IO \texttt{audio-chunk}} (server.north);
  \draw[arr] (server.south) -- node[lbl, right=2pt] {Persistent WebSocket} (ai.north);
  \draw[arr, dashed, draw=blue!50] (ai.east) -- ++(0.8,0) |- node[lbl, above, near start] {} (server.east);
  \draw[arr, dashed, draw=teal!50] (server.west) -- ++(-0.8,0) |- node[lbl, above, near start] {} (client.west);
  % Return labels
  \node[lbl, right=0.2cm] at ($(ai.east)+(0.85,0.5)$) {\texttt{text} + \texttt{audio\_chunk}};
  \node[lbl, left=0.2cm] at ($(server.west)+(-0.85,0.5)$) {\texttt{translation-result}};
\end{tikzpicture}%
}
\caption{Three-tier system architecture of SAMVADA. Solid arrows represent the primary data flow (audio upload). Dashed arrows represent response paths (translated text and synthesized audio).}
\label{fig:architecture}
\end{figure}

\subsection{Design Principles}

The system is designed around the following principles:

\begin{enumerate}
  \item \textbf{Subtitle-First Delivery:} Text translations are emitted to clients immediately after NMT completes, without waiting for TTS audio synthesis. This reduces perceived latency by 1--3 seconds.

  \item \textbf{Per-Speaker Isolation:} Each speaker has an independent processing queue. One speaker's slow audio processing does not block or delay another speaker's translations.

  \item \textbf{Graceful Degradation:} The system implements multi-level fallbacks at every stage---if Edge TTS fails, Sarvam AI is attempted; if that fails, gTTS is used. If the persistent WebSocket drops, a one-shot WebSocket fallback is used.

  \item \textbf{Stateless AI Service:} The FastAPI AI service is stateless and horizontally scalable. All session state is maintained in the Node.js server tier and MongoDB.

  \item \textbf{Zero Configuration for End Users:} Participants simply select their preferred language from a dropdown. All ASR language detection, NMT routing, and TTS voice selection happen automatically.
\end{enumerate}

% =============================================================================
% IV. AI PIPELINE COMPONENTS
% =============================================================================
\section{AI Pipeline Components}
\label{sec:pipeline}

\subsection{Speech-to-Text: Whisper ASR}

The ASR component uses \texttt{faster-whisper}~\cite{fasterwhisper}, a CTranslate2-based reimplementation of OpenAI Whisper that provides $4\times$ speedup through INT8 quantization.

\subsubsection{Model Selection and Quantization}

The Whisper-\texttt{small} variant (244M parameters) is selected as the optimal balance between Word Error Rate (WER) and inference latency for real-time applications. On an NVIDIA T4 GPU, the \texttt{INT8\_FLOAT16} compute type provides the best throughput.

\subsubsection{Voice Activity Detection}

Server-side VAD filtering is enabled with tuned parameters as shown in Table~\ref{tab:vad}.

\begin{table}[htbp]
\centering
\caption{Whisper VAD Filter Parameters}
\label{tab:vad}
\begin{tabular}{lrl}
\toprule
\textbf{Parameter} & \textbf{Value} & \textbf{Rationale} \\
\midrule
\texttt{threshold} & 0.55 & Aggressive noise filtering \\
\texttt{min\_speech\_dur.} & 300\,ms & Rejects short noise bursts \\
\texttt{min\_silence\_dur.} & 600\,ms & Faster sentence boundaries \\
\bottomrule
\end{tabular}
\end{table}

\subsubsection{Indic Language Prompt Conditioning}

Whisper's language model is primed with orthographically correct prompts for 9~Indic languages to prevent word-splitting artifacts. For example, without conditioning, Hindi text is incorrectly split by Whisper. The conditioning prompt primes the language model to output correct Devanagari orthography.

\subsubsection{Greedy Decoding with Repetition Penalty}

\begin{lstlisting}[language=Python, caption={Whisper transcription parameters}]
beam_size=1              # Greedy: 3-4x speedup
repetition_penalty=1.3   # Suppresses loops
suppress_blank=True      # Removes blank artifacts
no_speech_threshold=0.6  # Skips silent segments
hallucination_silence_threshold=2.0
\end{lstlisting}

Beam size~1 (greedy decoding) provides a 3--4$\times$ speedup over beam search on GPU with minimal quality degradation when combined with \texttt{repetition\_penalty=1.3}.

\subsection{Hallucination Filtering Framework}

A critical challenge in production ASR systems is \textbf{Whisper hallucination}---the model generates plausible but fabricated text when processing silence, background noise, or low-energy audio segments. SAMVADA implements a multi-stage filtering pipeline:

\textbf{Stage~1---Statistical Rejection:}
\begin{itemize}
  \item \texttt{no\_speech\_prob $>$ 0.60} $\rightarrow$ Reject (silence)
  \item \texttt{avg\_logprob $<$ --1.1} $\rightarrow$ Reject (low confidence)
  \item \texttt{lang\_prob $<$ 0.55} AND word count $\leq 6$ $\rightarrow$ Reject
  \item \texttt{lang\_prob $<$ 0.40} $\rightarrow$ Reject (critically low)
\end{itemize}

\textbf{Stage~2---Text-Based Filtering:}
\begin{itemize}
  \item Noise tag detection: \texttt{[MUSIC]}, \texttt{[APPLAUSE]}, etc.
  \item Filler noise words: ``um'', ``uh'', ``hmm'', ``ah'', ``oh''
  \item Explicit blocklist of $\sim$40 known hallucination phrases spanning 12~languages
  \item Fuzzy substring matching for multi-word spam phrases
  \item Repetition loop detection (same $n$-gram repeated $\geq 3$ times)
\end{itemize}

\textbf{Stage~3---Foreign Language Hallucination Detection:}
When Whisper processes faint breathing or background hiss, it often hallucinates short phrases in random languages. The foreign hallucination detector: (1)~checks if the detected language mismatches the user's expected language hint; (2)~rejects short utterances ($\leq 4$ words) with low confidence; and (3)~rejects all utterances with \texttt{lang\_prob $<$ 0.40}.

\subsection{Neural Machine Translation: NLLB-200}

The NMT component uses Meta's NLLB-200-distilled-600M model~\cite{nllb} converted to CTranslate2 INT8 format.

\subsubsection{Batched Translation}

For multi-participant rooms, translations to $N$ target languages are executed as a single batched \texttt{translate\_batch()} call, amortizing the encoder computation:

\begin{lstlisting}[language=Python, caption={Batched NMT translation}]
source_batch = [source_tokens] * len(targets)
target_prefixes = [[get_nllb_code(t)]
                   for t in targets]
results = translator.translate_batch(
    source_batch,
    target_prefix=target_prefixes,
    beam_size=effective_beam,
    max_decoding_length=128,
    repetition_penalty=1.1,
)
\end{lstlisting}

\subsubsection{Adaptive Beam Size}

Indic languages (Hindi, Tamil, Telugu, etc.) follow Subject-Object-Verb (SOV) word order, while English follows Subject-Verb-Object (SVO). Translating SOV $\rightarrow$ SVO requires significant structural reordering. SAMVADA uses \texttt{beam\_size=4} for Indic-to-English translation and \texttt{beam\_size=2} for all other pairs.

\subsubsection{Unicode Script-Based Language Resolution}

Whisper occasionally misclassifies the source language (e.g., detecting Devanagari text as Urdu). The \texttt{resolve\_source\_language()} function inspects Unicode code points to correct these misclassifications before invoking NLLB, as shown in Table~\ref{tab:unicode}.

\begin{table}[htbp]
\centering
\caption{Unicode Script-Based Language Resolution}
\label{tab:unicode}
\begin{tabular}{lll}
\toprule
\textbf{Unicode Range} & \textbf{Script} & \textbf{Language} \\
\midrule
U+0900--U+097F & Devanagari & Hindi \\
U+0980--U+09FF & Bengali & Bengali \\
U+0B80--U+0BFF & Tamil & Tamil \\
U+0C00--U+0C7F & Telugu & Telugu \\
U+0A80--U+0AFF & Gujarati & Gujarati \\
U+0C80--U+0CFF & Kannada & Kannada \\
U+0D00--U+0D7F & Malayalam & Malayalam \\
U+0A00--U+0A7F & Gurmukhi & Punjabi \\
\bottomrule
\end{tabular}
\end{table}

\subsection{Text-to-Speech: Multi-Engine Architecture}

SAMVADA implements a cascaded TTS engine selector with three tiers:

\begin{itemize}
  \item \textbf{Tier~1---Sarvam AI (Bulbul v3):} Sovereign Indian AI TTS supporting 10 regional Indian languages with native prosody. Activated when \texttt{USE\_SARVAM=true} and a valid API key is present.
  \item \textbf{Tier~2---Microsoft Edge Neural TTS:} Studio-grade neural voices for 35+ languages with zero cost. Uses named voices (e.g., \texttt{hi-IN-MadhurNeural}).
  \item \textbf{Tier~3---Google TTS (gTTS):} Universal fallback supporting 27 languages.
\end{itemize}

\textbf{Echo Prevention:} TTS synthesis is skipped when the target language matches the speaker's detected language, as the speaker's voice is already transmitted via WebRTC.

\textbf{Concurrent TTS:} All target language TTS calls are launched concurrently using \texttt{asyncio.gather()}, reducing multi-language TTS latency by approximately 50\%.

\subsection{CTranslate2 Thread Tuning}

The CTranslate2 runtime is configured with tuned thread parameters as shown in Table~\ref{tab:threads}.

\begin{table}[htbp]
\centering
\caption{CTranslate2 Thread Configuration for NLLB-200}
\label{tab:threads}
\begin{tabular}{lccc}
\toprule
\textbf{Parameter} & \textbf{GPU (T4)} & \textbf{CPU} & \textbf{Rationale} \\
\midrule
\texttt{inter\_threads} & 1 & 2 & GPU parallelism \\
\texttt{intra\_threads} & 4 & 4 & Optimal for cores \\
\bottomrule
\end{tabular}
\end{table}

% =============================================================================
% V. SYSTEM IMPLEMENTATION
% =============================================================================
\section{System Implementation}
\label{sec:implementation}

\subsection{Client Tier}

The client is built using \textbf{React~19} with \textbf{Vite~8} as the build tool.

\subsubsection{AudioWorklet-Based Voice Activity Detection}

Traditional approaches use \texttt{MediaRecorder} with fixed timeslice intervals to capture audio. SAMVADA replaces this with an \texttt{AudioWorkletNode} that processes audio frames every $\sim$20\,ms (960~samples at 48\,kHz), enabling fine-grained energy-based VAD.

The client-side VAD configuration is shown in Table~\ref{tab:clientvad}.

\begin{table}[htbp]
\centering
\caption{Client-Side AudioWorklet VAD Parameters}
\label{tab:clientvad}
\begin{tabular}{lrl}
\toprule
\textbf{Parameter} & \textbf{Value} & \textbf{Purpose} \\
\midrule
\texttt{SILENCE\_THRESHOLD} & 0.015 & RMS energy cutoff \\
\texttt{SILENCE\_DUR\_MS} & 800 & Silence before flush \\
\texttt{MIN\_SPEECH\_MS} & 300 & Min speech length \\
\texttt{MAX\_CHUNK\_MS} & 7500 & Hard cap per utterance \\
\texttt{FRAME\_MS} & 20 & Worklet frame size \\
\bottomrule
\end{tabular}
\end{table}

The VAD operates as follows:
\begin{itemize}
  \item \textbf{Speech Start:} When RMS energy exceeds \texttt{SILENCE\_THRESHOLD}, audio frames are accumulated.
  \item \textbf{Speech End:} After \texttt{SILENCE\_DURATION\_MS} of consecutive silence (and minimum \texttt{MIN\_SPEECH\_MS} of speech), accumulated frames are encoded as a self-contained WAV file (PCM Int16, 48\,kHz mono) and emitted via Socket.IO.
  \item \textbf{Hard Cap:} If speech continues for \texttt{MAX\_CHUNK\_MS} (7.5\,s), the chunk is force-flushed. The 7.5\,s cap is tuned for SOV languages (Hindi, Tamil) where the verb appears at the sentence end.
  \item \textbf{Leading Silence Suppression:} Pre-speech silence frames are not accumulated, preventing Whisper from receiving seconds of silence that inflate \texttt{no\_speech\_prob}.
\end{itemize}

\subsubsection{WAV Encoding}

Audio is encoded client-side into self-contained WAV files using a minimal 44-byte RIFF/WAV header implementation. This eliminates WebM EBML container overhead and provides a universally decodable format.

\subsubsection{Meeting User Interface}

The meeting interface comprises the components listed in Table~\ref{tab:components}.

\begin{table}[htbp]
\centering
\caption{Client-Side Meeting UI Components}
\label{tab:components}
\begin{tabular}{ll}
\toprule
\textbf{Component} & \textbf{Function} \\
\midrule
\texttt{PreJoinScreen} & Device selection, media preview \\
\texttt{MeetingRoom} & Main meeting layout \\
\texttt{ParticipantGrid} & Responsive video tile grid \\
\texttt{ParticipantTile} & Video tile + audio visualizer \\
\texttt{SubtitleOverlay} & Real-time translated subtitles \\
\texttt{TranscriptPanel} & Scrollable transcript sidebar \\
\texttt{LanguageSelector} & Mid-meeting language switcher \\
\texttt{MeetingToolbar} & Mic, camera, leave controls \\
\texttt{PipelineInspector} & AI diagnostic visualization \\
\texttt{MeetingHeader} & Room code, timer, participants \\
\bottomrule
\end{tabular}
\end{table}

\subsection{Server Tier}

The Node.js backend provides REST APIs and real-time Socket.IO communication.

\subsubsection{Translation Orchestrator}

The orchestrator manages the complete lifecycle of an audio chunk from receipt to translation delivery. Each speaker gets an independent FIFO queue (max depth~$= 4$). When the queue is full, the oldest chunk is evicted (freshest speech is prioritized). Chunks older than 8~seconds are silently dropped before processing.

This design ensures: (1)~one slow speaker never blocks another speaker's translations; (2)~only one AI job runs per speaker at a time (single-flight concurrency); and (3)~stale audio is automatically discarded.

\subsubsection{Intelligent Translation Routing}

The orchestrator computes the unique set of target languages needed by all participants in the room (excluding the speaker's own language), then makes a single AI service call for all required translations. Translation results are routed privately to each participant based on their \texttt{targetLanguage} preference.

\subsubsection{Persistent WebSocket AI Connection}

The \texttt{PersistentAIConnection} class maintains a single long-lived WebSocket to the FastAPI AI service. All audio processing requests are multiplexed over this connection using UUID-based \texttt{job\_id} tags. Table~\ref{tab:wscomp} compares the two connection architectures.

\begin{table}[htbp]
\centering
\caption{One-Shot vs.\ Persistent WebSocket Comparison}
\label{tab:wscomp}
\begin{tabular}{lcc}
\toprule
\textbf{Metric} & \textbf{One-Shot} & \textbf{Persistent} \\
\midrule
Conn.\ overhead & 80--200\,ms & 0\,ms \\
Requests before fail & 1 & Unlimited \\
Reconnection & Manual & Auto (exp.\ backoff) \\
Multiplexing & No & Yes (\texttt{job\_id}) \\
\bottomrule
\end{tabular}
\end{table}

The persistent connection auto-reconnects with exponential backoff (1\,s base, 30\,s cap). If unavailable during reconnection, the system transparently falls back to a one-shot WebSocket.

\subsubsection{Data Models}

The MongoDB schema uses three primary collections:
\begin{itemize}
  \item \textbf{User:} Authentication (local + Google OAuth), preferred language, avatar.
  \item \textbf{Meeting:} Room code, host, participants with per-participant \texttt{targetLanguage}, status lifecycle (\texttt{waiting} $\rightarrow$ \texttt{active} $\rightarrow$ \texttt{ended}), post-meeting AI summary.
  \item \textbf{Transcript:} Per-utterance records with speaker, source language, original text, and translations map.
\end{itemize}

\subsubsection{Security}
\begin{itemize}
  \item \textbf{Helmet.js:} HTTP security headers (CSP, HSTS, X-Frame-Options).
  \item \textbf{Rate Limiting:} Auth endpoints limited to 100~requests per 15-minute window.
  \item \textbf{JWT Authentication:} Token-based authentication with Socket.IO middleware verification.
  \item \textbf{Password Hashing:} bcrypt with salt rounds~$= 10$.
\end{itemize}

\subsection{AI Service Tier}

The FastAPI application exposes three interfaces:
\begin{enumerate}
  \item \textbf{REST Endpoint:} \texttt{POST /api/process-audio}---multipart form upload for batch processing.
  \item \textbf{WebSocket Endpoint:} \texttt{WS /ws/process-audio}---streaming bi-directional interface with \texttt{job\_id} multiplexing.
  \item \textbf{Interactive Dashboard:} \texttt{GET /demo}---self-contained HTML test interface.
\end{enumerate}

\subsubsection{GPU Warmup}

On startup, the lifespan context manager pre-loads both Whisper and NLLB models, then runs dummy inferences to pre-allocate CUDA kernel memory, JIT-compile lazy CUDA operations, and ensure the first real request is as fast as subsequent ones. Without warmup, the first request takes 2--5$\times$ longer.

\subsubsection{Streaming Pipeline}

The \texttt{SpeechToSpeechEngine.process\_stream()} is an async generator that yields results progressively:
\begin{enumerate}
  \item \texttt{\{type: ``text'', translations, diagnostics\}}---emitted immediately after NMT.
  \item \texttt{\{type: ``audio\_chunk'', lang, audio\_base64\}}---emitted per-language TTS.
  \item \texttt{\{type: ``done'', latency, diagnostics\}}---final metrics.
\end{enumerate}

This enables the subtitle-first delivery pattern: clients receive translated text typically 1--3~seconds before TTS audio arrives.

\subsection{Pipeline Diagnostic Engine}

SAMVADA includes a comprehensive diagnostic engine that analyzes all pipeline handoffs and pinpoints quality degradation, as summarized in Table~\ref{tab:diagnostics}.

\begin{table}[htbp]
\centering
\caption{Pipeline Diagnostic Checks and Severity Levels}
\label{tab:diagnostics}
\begin{tabular}{llc}
\toprule
\textbf{Stage} & \textbf{Diagnostic Check} & \textbf{Severity} \\
\midrule
VAD & Chunk $< 0.7$\,s duration & Warning \\
Language & Spoken $\neq$ hint language & Warning \\
STT & \texttt{avg\_logprob $<$ --0.95} & Warning \\
NMT & Empty or error translation & Error \\
NMT & Output length $\ll$ input & Warning \\
TTS & Fallback engine used & Warning \\
\bottomrule
\end{tabular}
\end{table}

% =============================================================================
% VI. EXPERIMENTAL EVALUATION
% =============================================================================
\section{Experimental Evaluation}
\label{sec:evaluation}

\subsection{Hardware Configuration}

Experiments were conducted on the hardware listed in Table~\ref{tab:hardware}.

\begin{table}[htbp]
\centering
\caption{Hardware Configuration}
\label{tab:hardware}
\begin{tabular}{ll}
\toprule
\textbf{Component} & \textbf{Specification} \\
\midrule
GPU (Cloud) & NVIDIA Tesla T4 (16\,GB VRAM) \\
CPU & AMD Ryzen 7 5800H (8C/16T) \\
RAM & 16\,GB DDR4 \\
GPU (Local) & NVIDIA RTX 3050 (4\,GB VRAM) \\
Network & ngrok / Cloudflare tunnel \\
\bottomrule
\end{tabular}
\end{table}

\subsection{Model Configurations}

The model configurations used are listed in Table~\ref{tab:models}.

\begin{table}[htbp]
\centering
\caption{Model Configurations Used in SAMVADA}
\label{tab:models}
\begin{tabular}{llll}
\toprule
\textbf{Comp.} & \textbf{Model} & \textbf{Params} & \textbf{Quant.} \\
\midrule
ASR & Whisper-small & 244M & INT8\_FP16 \\
NMT & NLLB-200-dist. & 600M & INT8 \\
TTS & Edge Neural & --- & --- \\
TTS & Sarvam v3 & --- & API \\
\bottomrule
\end{tabular}
\end{table}

\subsection{Translation Quality Benchmark}

A 10-sentence English-to-Hindi benchmark was used to evaluate translation quality using SacreBLEU~\cite{sacrebleu}. Sample sentences are shown in Table~\ref{tab:benchmark}.

\begin{table}[htbp]
\centering
\caption{Sample Sentences from the EN$\rightarrow$HI Benchmark}
\label{tab:benchmark}
\begin{tabular}{cL{3.3cm}L{3.3cm}}
\toprule
\textbf{\#} & \textbf{English Source} & \textbf{Hindi Reference} \\
\midrule
1 & Hello, how are you today? & Namaste, aaj aap kaise hain? \\
2 & The AI system is processing the audio. & AI pranali audio ko process kar rahi hai. \\
3 & We need to schedule a meeting for tomorrow. & Humein kal ke liye ek baithak tay karni hogi. \\
\bottomrule
\end{tabular}
\end{table}

\subsection{Latency Measurements}

End-to-end latency is decomposed into four stages:
\begin{equation}
T_{\text{total}} = T_{\text{VAD}} + T_{\text{ASR}} + T_{\text{NMT}} + T_{\text{TTS}}
\end{equation}

Where $T_{\text{VAD}}$ is the client-side voice activity detection and WAV encoding ($\sim$800\,ms silence wait), $T_{\text{ASR}}$ is Whisper speech-to-text inference, $T_{\text{NMT}}$ is NLLB translation to $N$ target languages, and $T_{\text{TTS}}$ is Edge Neural TTS synthesis.

\subsection{Supported Languages}

SAMVADA supports 30~languages across 10~language families. Table~\ref{tab:languages} summarizes the language coverage.

\begin{table}[htbp]
\centering
\caption{Supported Languages Grouped by Family}
\label{tab:languages}
\renewcommand{\arraystretch}{1.15}
{\small
\begin{tabular}{ll}
\toprule
\textbf{Language Family} & \textbf{Languages Supported} \\
\midrule
Indo-European (Germanic) & English, German, Dutch \\
Indo-European (Romance) & French, Spanish, Portuguese, Italian \\
Indo-European (Slavic) & Russian, Polish, Ukrainian \\
Indo-Aryan & Hindi, Marathi, Bengali, Gujarati, \\
 & Punjabi, Urdu, Odia, Assamese \\
Dravidian & Tamil, Telugu, Kannada, Malayalam \\
Turkic & Turkish \\
Sino-Tibetan & Chinese (Simplified) \\
Other & Japanese, Korean, Arabic, \\
 & Swahili, Thai, Vietnamese \\
\bottomrule
\end{tabular}
}
\end{table}

% =============================================================================
% VII. RESULTS AND DISCUSSION
% =============================================================================
\section{Results and Discussion}
\label{sec:results}

\subsection{Latency Analysis}

The measured average latencies for the pipeline stages on an NVIDIA T4 GPU are shown in Table~\ref{tab:latency}.

\begin{table}[htbp]
\centering
\caption{Average Pipeline Latency Measurements}
\label{tab:latency}
\begin{tabular}{lcc}
\toprule
\textbf{Stage} & \textbf{GPU} & \textbf{CPU} \\
\midrule
ASR (beam=1) & 0.3--0.5\,s & 1.5--2.5\,s \\
NMT (1 target) & 0.08--0.15\,s & 0.3--0.5\,s \\
NMT (3 targets) & 0.12--0.20\,s & 0.5--0.8\,s \\
TTS (1 lang) & 0.5--1.5\,s & 0.5--1.5\,s \\
\midrule
\textbf{Subtitle} & \textbf{0.4--0.7\,s} & \textbf{1.8--3.0\,s} \\
\textbf{w/ Audio} & \textbf{1.5--3.5\,s} & \textbf{3.0--5.0\,s} \\
\bottomrule
\end{tabular}
\end{table}

Key observations:
\begin{enumerate}
  \item \textbf{Subtitle-first delivery} provides text translations in under 1~second on GPU, which is below the human perception threshold for real-time subtitles.
  \item \textbf{Greedy decoding} (\texttt{beam\_size=1}) reduces Whisper ASR latency by 3--4$\times$ compared to \texttt{beam\_size=5}.
  \item \textbf{Batched NMT} for 3~languages adds only $\sim$40--60\,ms overhead compared to single-language translation.
  \item \textbf{TTS is the bottleneck}---Edge Neural TTS is network-dependent. Concurrent \texttt{asyncio.gather()} reduces multi-language TTS from $\mathcal{O}(N)$ to $\mathcal{O}(1)$ in wall-clock time.
\end{enumerate}

\subsection{Connection Architecture Impact}

The persistent WebSocket eliminates per-request connection overhead of 80--200\,ms (TCP + TLS + ngrok handshake). This represents a 100\% reduction in connection overhead, unlimited request throughput via multiplexing, and automated reconnection with exponential backoff.

\subsection{Hallucination Filtering Effectiveness}

Without filtering, approximately 15--25\% of processed audio chunks in a typical meeting contain Whisper hallucinations (primarily during pauses between speakers). The multi-stage filtering framework reduces this to $<$2\%, catching silence artifacts, YouTube/streaming credit hallucinations (12~languages), repetition loops, wrong-language short phrases, and noise tag markers.

\subsection{Limitations}

\begin{enumerate}
  \item \textbf{NLLB-200 Quality for Indic Languages:} While NLLB-200 supports 200+ languages, its translation quality for low-resource Indic language pairs is lower than for high-resource pairs. Specialized models like IndicTrans2~\cite{indictrans2} offer superior quality but at higher computational cost.

  \item \textbf{No End-to-End Model:} The cascaded pipeline introduces error propagation---ASR errors compound in NMT, which compounds in TTS.

  \item \textbf{TTS Naturalness:} Edge Neural TTS produces high-quality speech but does not retain the original speaker's voice. Voice cloning using XTTS-v2~\cite{xtts} is architecturally supported but not yet activated due to VRAM constraints.

  \item \textbf{Network Dependency:} TTS engines require internet connectivity. In offline environments, TTS would fail, falling back to text-only subtitles.

  \item \textbf{Evaluation Scale:} The benchmark dataset contains only 10~sentences. A more comprehensive evaluation using standardized datasets (e.g., FLORES-200~\cite{nllb}) would provide statistically significant quality metrics.
\end{enumerate}

% =============================================================================
% VIII. CONCLUSION AND FUTURE WORK
% =============================================================================
\section{Conclusion and Future Work}
\label{sec:conclusion}

\subsection{Conclusion}

This paper presented SAMVADA, a real-time Speech-to-Speech Translation platform for multilingual video conferencing. The system successfully demonstrates that commodity GPU hardware combined with optimized open-source models (Whisper, NLLB-200) and cloud TTS services can deliver sub-4-second end-to-end S2ST for 30~languages. Key technical contributions include the per-speaker priority queue orchestrator, persistent multiplexed WebSocket architecture, multi-stage Whisper hallucination filtering, AudioWorklet-based continuous VAD, and the pipeline diagnostic engine.

The platform is implemented as a full-stack web application with a React frontend, Node.js backend, and FastAPI AI service, making it accessible via any modern web browser without additional software installation. The modular architecture allows individual pipeline components to be upgraded independently as better models become available.

\subsection{Future Work}

\begin{enumerate}
  \item \textbf{Voice Cloning Integration:} Activate the XTTS-v2 voice retention engine to synthesize translated speech in the original speaker's voice.

  \item \textbf{End-to-End Models:} Evaluate Meta's SeamlessM4T as a drop-in replacement for the cascaded pipeline.

  \item \textbf{IndicTrans2 Hybrid Router:} Implement a language-aware routing layer using AI4Bharat's IndicTrans2 for Indian language pairs and NLLB-200 for all other pairs.

  \item \textbf{On-Device ASR:} Explore Whisper.cpp and WebAssembly-based ASR to perform speech recognition directly in the browser.

  \item \textbf{Comprehensive Evaluation:} Conduct systematic evaluation using FLORES-200, CoVoST-2, and FLEURS benchmarks across all 30~supported languages.

  \item \textbf{Speaker Diarization:} Integrate speaker diarization to handle overlapping speech more accurately.

  \item \textbf{Streaming ASR:} Replace VAD-chunked batch ASR with true streaming ASR (e.g., Whisper-Streaming) to reduce $T_{\text{VAD}}$.
\end{enumerate}

% =============================================================================
% ACKNOWLEDGMENTS
% =============================================================================
\section*{Acknowledgments}

The authors acknowledge the open-source contributions of OpenAI (Whisper), Meta~AI (NLLB-200), Microsoft (Edge Neural TTS), Sarvam~AI (Bulbul TTS), the faster-whisper project, and the CTranslate2 team. The Google Colab platform provided GPU compute resources for development and testing.

\balance

% =============================================================================
% REFERENCES
% =============================================================================
\begin{thebibliography}{15}

\bibitem{ethnologue}
D.~M. Eberhard, G.~F. Simons, and C.~D. Fennig, Eds., \emph{Ethnologue: Languages of the World}, 27th~ed. Dallas, TX: SIL International, 2024. [Online]. Available: \url{https://www.ethnologue.com}

\bibitem{nakamura2006atr}
S.~Nakamura \emph{et~al.}, ``The ATR multilingual speech-to-speech translation system,'' \emph{IEEE Trans. Audio, Speech, Lang. Process.}, vol.~14, no.~2, pp.~365--376, Mar.~2006.

\bibitem{seamlessm4t}
Meta~AI, ``SeamlessM4T: Massively multilingual \& multimodal machine translation,'' \emph{arXiv preprint arXiv:2308.11596}, 2023.

\bibitem{translatotron2}
Y.~Jia \emph{et~al.}, ``Translatotron~2: High-quality direct speech-to-speech translation with voice preservation,'' in \emph{Proc. ICML}, 2022.

\bibitem{whisper}
A.~Radford \emph{et~al.}, ``Robust speech recognition via large-scale weak supervision,'' in \emph{Proc. ICML}, 2023.

\bibitem{fasterwhisper}
SYSTRAN, ``faster-whisper: Faster Whisper transcription with CTranslate2,'' GitHub, 2023. [Online]. Available: \url{https://github.com/SYSTRAN/faster-whisper}

\bibitem{nllb}
NLLB~Team \emph{et~al.}, ``No language left behind: Scaling human-centered machine translation,'' \emph{arXiv preprint arXiv:2207.04672}, 2022.

\bibitem{edgetts}
Microsoft, ``Microsoft Edge text-to-speech,'' 2024. [Online]. Available: \url{https://azure.microsoft.com/en-us/services/cognitive-services/text-to-speech/}

\bibitem{sarvam}
Sarvam~AI, ``Bulbul v3: Indian language TTS model,'' 2024. [Online]. Available: \url{https://www.sarvam.ai/}

\bibitem{gtts}
P.~N. Suganthan, ``gTTS: Google text-to-speech,'' GitHub, 2023. [Online]. Available: \url{https://github.com/pndurette/gTTS}

\bibitem{webrtc}
A.~B. Johnston and D.~C. Burnett, \emph{WebRTC: APIs and RTCWEB Protocols of the HTML5 Real-Time Web}, 3rd~ed. Digital Codex LLC, 2014.

\bibitem{socketio}
Socket.IO, ``Socket.IO: Bidirectional and low-latency communication,'' 2024. [Online]. Available: \url{https://socket.io/}

\bibitem{sacrebleu}
M.~Post, ``A call for clarity in reporting BLEU scores,'' in \emph{Proc. WMT}, 2018.

\bibitem{indictrans2}
AI4Bharat, ``IndicTrans2: Towards high-quality and accessible machine translation models for all 22 scheduled Indian languages,'' \emph{Trans. Mach. Learn. Res.}, 2023.

\bibitem{xtts}
Coqui~AI, ``XTTS-v2: Cross-lingual text-to-speech,'' GitHub, 2023. [Online]. Available: \url{https://github.com/coqui-ai/TTS}

\end{thebibliography}

% =============================================================================
% APPENDICES
% =============================================================================
\appendices

\section{Technology Stack Summary}

Table~\ref{tab:techstack} lists the complete technology stack used across all three tiers of the SAMVADA platform.

\begin{table*}[htbp]
\centering
\caption{Complete Technology Stack of SAMVADA}
\label{tab:techstack}
\renewcommand{\arraystretch}{1.2}
{\small
\begin{tabular}{@{} l l l l l @{}}
\toprule
\textbf{Tier} & \textbf{Technology} & \textbf{Version} & \textbf{Purpose} & \textbf{License} \\
\midrule
\multirow{4}{*}{\textbf{Frontend}} 
  & React          & 19.2.7   & UI component framework              & MIT \\
  & Vite           & 8.1.1    & Build tool \& dev server             & MIT \\
  & Socket.IO Client & 4.8.3  & Real-time bi-directional transport   & MIT \\
  & React Router   & 7.18.1   & Client-side routing                  & MIT \\
\midrule
\multirow{5}{*}{\textbf{Backend}}
  & Express        & 4.19.2   & REST API framework                   & MIT \\
  & Socket.IO      & 4.7.5    & WebSocket server                     & MIT \\
  & Mongoose       & 8.4.1    & MongoDB object data modeling          & MIT \\
  & JSON Web Token & 9.0.2    & Stateless authentication             & MIT \\
  & Helmet         & 8.3.0    & HTTP security headers                & MIT \\
\midrule
\multirow{7}{*}{\textbf{AI Service}}
  & FastAPI        & 0.111.0  & Async ML API server                  & MIT \\
  & PyTorch        & Latest   & Deep learning framework              & BSD-3 \\
  & faster-whisper & $\geq$1.0.0 & ASR engine (CTranslate2 backend)  & MIT \\
  & CTranslate2    & Latest   & Optimized inference runtime           & MIT \\
  & Transformers   & 4.44.2   & NLLB tokenizer                       & Apache-2.0 \\
  & edge-tts       & $\geq$6.1.9 & Microsoft Edge Neural TTS         & MIT \\
  & gTTS           & 2.5.1    & Google Text-to-Speech fallback       & MIT \\
\bottomrule
\end{tabular}
}
\end{table*}

\section{API Endpoint Reference}

Table~\ref{tab:endpoints} documents all REST and WebSocket API endpoints exposed by the server and AI service tiers.

\begin{table*}[htbp]
\centering
\caption{REST and WebSocket API Endpoints}
\label{tab:endpoints}
\renewcommand{\arraystretch}{1.2}
{\small
\begin{tabular}{@{} l l l l @{}}
\toprule
\textbf{Method} & \textbf{Endpoint} & \textbf{Service} & \textbf{Description} \\
\midrule
\multicolumn{4}{l}{\textit{Server Tier --- Node.js Backend}} \\
\midrule
GET  & \texttt{/api/health}               & Server  & Server health check and status \\
POST & \texttt{/api/auth/register}         & Server  & New user registration (email + password) \\
POST & \texttt{/api/auth/login}            & Server  & User login with JWT token issuance \\
POST & \texttt{/api/auth/google}           & Server  & Google OAuth 2.0 authentication \\
POST & \texttt{/api/meetings}              & Server  & Create a new meeting room \\
GET  & \texttt{/api/meetings/:id}          & Server  & Retrieve meeting details and participants \\
POST & \texttt{/api/meetings/:id/summarize}& Server  & Generate AI-powered meeting summary \\
GET  & \texttt{/api/meetings/:id/summary}  & Server  & Retrieve stored meeting summary \\
GET  & \texttt{/api/transcripts/:meetingId}& Server  & Retrieve all transcripts for a meeting \\
\midrule
\multicolumn{4}{l}{\textit{AI Service Tier --- FastAPI Backend}} \\
\midrule
GET  & \texttt{/health}                    & AI      & AI service health with loaded model info \\
GET  & \texttt{/api/languages}             & AI      & Supported language codes and display names \\
POST & \texttt{/api/process-audio}         & AI      & Full S2ST pipeline via multipart upload \\
WS   & \texttt{/ws/process-audio}          & AI      & Real-time streaming S2ST via WebSocket \\
GET  & \texttt{/demo}                      & AI      & Interactive test dashboard (HTML) \\
\bottomrule
\end{tabular}
}
\end{table*}

\section{Socket.IO Event Protocol}

Table~\ref{tab:socketevents} specifies the complete bi-directional Socket.IO event protocol used for real-time communication between clients and the server.

\begin{table*}[htbp]
\centering
\caption{Socket.IO Event Protocol Specification}
\label{tab:socketevents}
\renewcommand{\arraystretch}{1.2}
{\small
\begin{tabular}{@{} c l l l @{}}
\toprule
\textbf{Direction} & \textbf{Event Name} & \textbf{Key Payload Fields} & \textbf{Description} \\
\midrule
\multicolumn{4}{l}{\textit{Client $\to$ Server Events}} \\
\midrule
C $\to$ S & \texttt{join-meeting}       & \texttt{roomCode, userId, displayName, targetLanguage} & Join a meeting room \\
C $\to$ S & \texttt{audio-chunk}        & \texttt{ArrayBuffer + \{roomCode, speakerName, mimeType\}} & Send captured audio for translation \\
C $\to$ S & \texttt{update-language}     & \texttt{roomCode, targetLanguage} & Change preferred receive language \\
C $\to$ S & \texttt{leave-meeting}       & \texttt{roomCode, userId} & Leave the current meeting \\
C $\to$ S & \texttt{end-meeting}         & \texttt{roomCode} & Host terminates meeting for all \\
C $\to$ S & \texttt{toggle-media}        & \texttt{roomCode, userId, isMuted, isVideoOff} & Broadcast mic/camera toggle \\
\midrule
\multicolumn{4}{l}{\textit{Server $\to$ Client Events}} \\
\midrule
S $\to$ C & \texttt{meeting-joined}      & \texttt{participants[ ]} & Confirm join with participant list \\
S $\to$ C & \texttt{translation-pending}  & \texttt{speakerId, speakerName, timestamp} & Processing indicator (``Translating\ldots'') \\
S $\to$ C & \texttt{translation-result}   & \texttt{originalText, translatedText, lang, diagnostics} & Translated subtitle text \\
S $\to$ C & \texttt{translation-audio}    & \texttt{Buffer + \{mimeType, lang, speakerName\}} & Synthesized TTS audio binary \\
S $\to$ C & \texttt{new-transcript}       & \texttt{speakerName, originalText, translations, sequenceNumber} & Transcript entry for sidebar \\
S $\to$ C & \texttt{speaker-subtitle}     & \texttt{originalText, isSelf: true} & Speaker's own STT feedback \\
S $\to$ C & \texttt{translation-diagnostics} & \texttt{diagnostics \{status, warnings, stt, nmt, tts\}} & Pipeline health report \\
S $\to$ C & \texttt{participant-joined}   & \texttt{userId, displayName, targetLanguage} & New participant notification \\
S $\to$ C & \texttt{participant-left}     & \texttt{userId, socketId} & Participant departure notification \\
S $\to$ C & \texttt{meeting-ended}        & \texttt{message, roomCode} & Host ended the meeting \\
\bottomrule
\end{tabular}
}
\end{table*}

\end{document}
